\section{Introduction}
\label{sec:introduction}

H5P is an open-source framework for interactive learning content.
An H5P content item is a
self-contained package that a learning platform embeds into its pages, and the Moodle learning
management system supports it out of the box: a teacher uploads a package, and students read it in
the browser as part of the course, with nothing to install and no external site to visit.
That
makes H5P attractive for course material of every kind, from a page of text to a set of
interactive questions.
Creating H5P content, however, means working in the platform's browser-based
editor.
Writing more than a few pages this way is slow, and
revising what has been written is slower still.

Markdown, on the other hand, is plain text with a handful of conventions for headings,
emphasis, lists, links, and code, so a chapter can be written in any text editor, kept in a Git
repository, reviewed as a diff, and shared without needing a particular program.
What Markdown lacks is a comfortable way to be read inside a course: a \code{.md} file attached to a Moodle page
is a download, and a student who opens it sees the raw syntax unless they have a viewer.
What is convenient to write is therefore not what is convenient to read.
The usual answer is to port the material by hand: open the H5P editor, paste and reformat each paragraph, upload
each image, rebuild each table.
This is tedious and must be repeated whenever the Markdown source changes to keep the source and the H5P copy in sync.

Among H5P's content types, Interactive Book is the one that best fits
a folder of Markdown files: it presents content as a sequence of chapters with a table of
contents, and each chapter is a list of content blocks such as text, images, and tables.
A folder
of Markdown files, one per chapter, maps onto it almost directly.

Almost, because the two formats do not line up completely, in both directions.
H5P has content
types for which Markdown has no syntax, such as quizzes or image hotspots.
Markdown, on the other
hand, permits structures that H5P refuses: an image inside a table cell is ordinary Markdown, but H5P table blocks may not contain images.
A converter therefore has to make compromises. The subject of this thesis is
which compromises are possible, which were chosen, and why.

The thesis makes five contributions.
(1) It works out what Moodle's H5P implementation actually
accepts in a text block, which HTML tags and which CSS declarations, by reading the sanitizer and
the content-type definitions, since the documentation does not say.
(2) For every
Markdown construct that has no direct counterpart, it presents a rendering strategy.
(3) It presents md2h5p, a
converter that turns a folder of Markdown files, a Git repository, or an upload into an Interactive
Book package and reports constructs it could only approximate as warnings.
(4) It adds a
small Markdown extension for single- and multiple-choice questions that leaves the source file
valid Markdown for every other tool.
(5) It adds a merge step that regenerates a book from changed
Markdown without losing what was added to the previous version in the H5P editor.
